\documentclass[fleqn,a4paper]{article} % Flush Left Equation Numbers

% --- PACKAGES ---
\usepackage{graphicx} % Required for inserting images
\usepackage{a4wide} % Tell LaTeX to use the room there is on the page
\usepackage[calc,showdow,english]{datetime2} % To get the date of today
\usepackage{amsmath} % For matrices
\usepackage{algorithm} % Algorithm
\usepackage{algpseudocode} % Algorithm
\usepackage{subcaption} % For captions of subfigures
\usepackage{booktabs} % toprule
\usepackage[hidelinks,colorlinks=true,citecolor=red,urlcolor=blue]{hyperref} % Make clickable references which are less ugly (borders removed)
\usepackage[style=phys, articletitle=false,biblabel=brackets, chaptertitle=false,pageranges=false]{biblatex} % APS style (American Physical Society)
\usepackage[absolute]{textpos}
\usepackage{bm} % For bold Greek letters
\usepackage{siunitx} % For units in equations
\numberwithin{equation}{section} % For equation numbering
\numberwithin{figure}{section} % For figure numbering
\numberwithin{algorithm}{section} % For algorithm numbering

\addbibresource{references.bib}

\begin{document}
\section{Continuous equations}
We have a two-component fluid. The continuity equations are given as follows,
\begin{equation}
    \frac{\partial \rho_i}{\partial t} + \bm{\nabla}\cdot \Bigl(\rho_i \bm{u}_i\Bigr) = 0. 
\end{equation}
We can define a composition field $\phi = \rho_1 - \rho_2$,
\begin{equation}\label{eq:composition field}
    \frac{\partial \phi}{\partial t} + \bm{\nabla}\cdot \Bigl(\phi \bm{u}\Bigr) = - \bm{\nabla} \cdot \left[\frac{\rho}{2}\left(1 - \frac{\phi^2}{\rho^2}\right)\left(\bm{u}_1 - \bm{u}_2\right)\right].
\end{equation}
The momentum equation for the total density can be given as follows,
\begin{equation}
    \frac{\partial \left(\rho \bm{u}\right)}{\partial t} + \bm{\nabla} \cdot\Bigl(\rho \bm{u}\bm{u}\Bigr) = \bm{\nabla}\cdot \underline{\underline{\tau}} + \bm{\nabla}\cdot \underline{\underline{P}} + \bm{f}_b. 
\end{equation}
The viscous stress tensor was denoted by $\underline{\underline{\tau}}$, 
\begin{equation}
    \tau_{\alpha \beta} = \mu \left[\partial_\alpha u_\beta + \partial_\beta u_\alpha - \frac23 \partial_\mu u_\mu \delta_{\alpha \beta}\right] + \zeta \partial_\mu u_\mu \delta_{\alpha \beta}.
\end{equation}
The pressure tensor was denoted by $\underline{\underline{P}}$, and includes for instance the hydrodynamic pressure and the surface tension. Lastly, $\bm{f}_b$ denotes arbitrary body forces on the fluid.
\section{Free Energy}
A nonideal fluid can be simulated using the Free Energy Lattice Boltzmann Model~\cite{swift_lattice_1996}. We start by defining the Free Energy functional $\Psi$. This functional will depend on the fluid density $\rho$ and its gradient $\bm{\nabla}\rho$. An example for a one-component fluid can be given as follows,
\begin{equation}
    \Psi[\rho, \bm{\nabla}\rho] = \int \left(\psi_b + \frac{\kappa}{2}|\bm{\nabla}\rho|^2\right)\mathrm{d}V,
\end{equation}
where $\kappa$ is related to the surface tension and $\psi_b$ is the bulk free energy. An example for the bulk free energy can be given by the Van der Waals free energy~\cite{swift_lattice_1996},
\begin{equation}
    \psi_b = \rho T \ln \left(\frac{\rho}{1-\rho b}\right) - a \rho^2. 
\end{equation}
We can define the chemical potential as follows,
\begin{equation}
    \mu \equiv \frac{\delta \Psi}{\delta \rho}. 
\end{equation}
When we have a binary fluid, we will define separate chemical potentials for the total density $\rho$ and the composition $\phi = \rho_1 - \rho_2$,
\begin{equation}
    \mu_\rho \equiv \frac{\delta \Psi}{\delta \rho},\quad \mu_\phi \equiv \frac{\delta \Psi}{\delta \phi}.
\end{equation}
We can find a relation between the pressure tensor and the chemical potentials,
\begin{equation}
    \partial_\beta P_{\alpha \beta} = \rho_1 \partial_\alpha \mu_{\rho_1} + \rho_2 \partial_\alpha \mu_{\rho_2}
\end{equation}
I will use the functional chainrule,
\begin{equation}
    \frac{\delta \Psi}{\delta \rho_i} = \frac{\delta \Psi}{\delta \rho}\frac{\partial \rho}{\partial \rho_i} + \frac{\delta \Psi}{\delta \phi}\frac{\partial \phi}{\partial \rho_i}.
\end{equation}
This gives the pressure tensor in terms of functional derivatives over the total density and composition,
\begin{equation}\label{eq:pressure tensor}
    \partial_\beta P_{\alpha \beta} = \rho \partial_\alpha \mu_\rho + \phi \partial_\alpha \mu_\phi. 
\end{equation}
\section{Composition field complications}
The literature is divided on how the composition field is modeled. Some literature solved Eq.~(\ref{eq:composition field}) directly~\cite{swift_lattice_1996}\cite{kruger_lattice_2017}. However, they model a modified equation that captures the velocity difference using the gradient of a chemical potential,
\begin{equation}\label{eq:cahn-hilliard}
    \frac{\partial \phi}{\partial t} + \bm{\nabla}\cdot \Bigl(\phi \bm{u}\Bigr) = \bm{\nabla} \cdot \Bigl[M\bm{\nabla}\mu_{\phi}\Bigr].
\end{equation}
One issue is that the foundational paper gets an additional term added to this equation~\cite{swift_lattice_1996},
\begin{equation}\label{eq:cahn-hilliard added}
    \frac{\partial \phi}{\partial t} + \bm{\nabla}\cdot \Bigl(\phi \bm{u}\Bigr) = \bm{\nabla} \cdot \Bigl[M\bm{\nabla}\mu_{\phi}\Bigr] - \bm{\nabla}\cdot \left(\frac{\phi}{\rho}\bm{\nabla}\cdot \underline{\underline{P}}\right).
\end{equation}
I am confused where this term comes from. Later literature uses the exact same computational method, but communicates Eq.~(\ref{eq:cahn-hilliard}) as resulting equation from the numerical scheme, instead of Eq.~(\ref{eq:cahn-hilliard added}). I will assume the newer literature is correct. 

\section{Numerical scheme}
I will apply the mixed Free Energy Model~\cite{kruger_lattice_2017}. We rewrite Eq.~(\ref{eq:pressure tensor}),
\begin{equation}\label{eq:pressure tensor rewritten}
    \partial_\beta P_{\alpha \beta} = \Bigl[\rho \partial_\alpha \mu_\rho + \partial_\alpha \left(\phi \mu_\phi\right)\Bigr] - \mu_\phi \partial_\alpha \phi
\end{equation}
Usually, this can be rewritten in the following form,
\begin{equation}
    \partial_\beta P_{\alpha \beta} = \Bigl[\partial_\alpha \left(\rho / 3\right) + \partial_\alpha \left(\phi \mu_\phi\right)\Bigr] - \mu_\phi \partial_\alpha \phi
\end{equation}
The first part on the right hand side will be added to the equilibrium distribution. We will solve the total fluid motion with a distribution function $f$, 
\begin{equation}
    \rho = \sum_i f_i, \quad \rho \bm{u} = \sum_i \bm{c}_i f_i + \frac12 \bm{F}, 
\end{equation}
where $\bm{F}$ is additional body forces, and $\bm{c}_i$ are the discrete velocities. The equilibrium distribution is given as follows,
\begin{equation}
    f_i^\text{eq} = w_i \rho \left(1 + 3\frac{\phi \mu_\phi}{\rho} + 3 \bm{u} \cdot \bm{c}_i + \frac{9}{2}\left(\bm{u} \cdot \bm{c}_i\right)^2 - \frac32|\bm{u}|^2\right).
\end{equation}
Note that we used $f_0^\text{eq} = \rho - \sum_{i\neq 0} f_i^\text{eq}$. The second term on the right hand side in Eq.~(\ref{eq:pressure tensor rewritten}) will be modeled by an additional force,
\begin{equation}
    \bm{F} = -\mu_\phi \bm{\nabla}\phi. 
\end{equation}
Finally, we will also solve Eq.~(\ref{eq:cahn-hilliard}) using a distribution function $g$, with $\phi = \sum_i g_i$. We use the following equilibrium,
\begin{equation}
    g_i^\text{eq} = w_i \left(3\Gamma \mu_\phi + 3 \phi \bm{u}\cdot \bm{c}_i + \frac{9}{2}\phi\left(\bm{u} \cdot \bm{c}_i\right)^2 -  \frac32\phi|\bm{u}|^2\right).
\end{equation}
Note that we used $g_0^\text{eq} = \phi - \sum_{i\neq 0}g_i^\text{eq}$. Finally, $M = \Gamma \left(\tau_\phi - \frac12\right)$, where $\tau_\phi$ is the relaxation time for the Lattice Boltzmann algorithm. 

\clearpage % Restrict floats (images)
\printbibliography{}

\end{document}